% id: 203-23
% book: 베이직쎈 대수 · 12 수학적 귀납법
% section: 유형 127 수학적 귀납법; 부등식의 증명
% figure: none
% answer: 
% answer_source: 
% variant_level: 0 (원본 전사)
% 이 파일은 build.mjs 가 items.json 에서 만든다. 고칠 때는 items.json 을 고친다.
\smash{\raisebox{1.5mm}{\dmkichul{베이직쎈 대수 203쪽 23번}}}
다음은 $n\ge 5$인 모든 자연수 $n$에 대하여 부등식 $2^n>n^2$이 성립함을 수학적 귀납법으로 증명한 것이다.\exprbox{\begin{minipage}{0.96\linewidth}\raggedright\lineskip=4pt {\small\bfseries 증명}\par\smallskip (i) $n=\fbox{㈎}$일 때,\\ \hspace*{3em}(좌변)$=2^5=32$, (우변)$=5^2=25$\\ \hspace*{1.5em}따라서 주어진 부등식이 성립한다.\\ (ii) $n=k$ $(k\ge 5)$일 때 주어진 부등식이 성립한다고 가정하면 \quad $2^k>k^2$\\ \hspace*{1.5em}위의 식의 양변에 $2$를 곱하면\\ \hspace*{3em}$2^{k+1}>2k^2$ \quad $\cdots\cdots$ ㉠\\ \hspace*{1.5em}그런데 $k\ge 5$이면 $k^2-2k-1=\fbox{㈏}-2>0$이므로\\ \hspace*{3em}$k^2>2k+1$ \quad $\cdots\cdots$ ㉡\\ \hspace*{1.5em}㉠, ㉡에서 \quad $2^{k+1}>2k^2=k^2+k^2>\fbox{㈐}$\\ \hspace*{1.5em}따라서 $n=k+1$일 때도 주어진 부등식이 성립한다.\\ (i), (ii)에서 $n\ge 5$인 모든 자연수 $n$에 대하여 주어진 부등식이 성립한다.\end{minipage}}위의 과정에서 ㈎, ㈏, ㈐에 알맞은 것을 구하시오.
